% !TeX root = ../papers/nc-p3.tex
\section{階数15曲線族の無限遠八点の重み付き極限}
\label{sec:rank-fifteen-pencil-infinity}
\begin{theorem}\label{thm:rank-fifteen-pencil-infinity}
階数15の各表示で、一般の水準 $\xi=c^2$ の $K=\infty$ 上の八点は
全て、矢の基底を取り直すと標準Koszul左右AS正則な極限を持つ。
これは二つの種数3の特殊水準でも全八点で成り立つ。さらにこの二水準では、
109上のある付値で $K$ が極を持つ任意の通常な標数零窓が正則であり、
各始点で四次中心の多項式環上の階数256の自由加群となる。
\end{theorem}
\begin{proof}
$u=1/K$、$L^2=1+c^2$ とし、$un$ の定数項が $1/L$ となるように $L$ の符号を選ぶ。
$v=L+c$ により $L=(v+v^{-1})/2$、$c=(v-v^{-1})/2$ となる。
$S=2c/L^2$ と置く。前節の二次式から
$r,w$ の各根の先頭項は、それぞれ
\[
 r\sim S u^{-2}\text{ または }u^2/S,\qquad
 w\sim -S u^{-2}\text{ または }-u^2/S
\]
となる。これらは指数を除けば単元の級数で、逆数は小さい方の根を与える。
例えば $u^2r$ について大きな根の方を解くと定数項 $S$ は単根である。
$n$ の二根も $un$ について単根となる。$L$ の二符号と $r,w$ の二選択で八点を尽くす。

$r\sim R u^d$、$w\sim W u^e$ と書く。第一表示の $a,b,e',f$（$e'$ は窓の第三座標）の
先頭項は順に
\[
 L^{-1}u^{-1},\quad (1+c/L)R^{-1}u^{-1-d},\quad
 -L^{-1}u^{-1},\quad(1-c/L)R^{-1}u^{-1-d}.
\]
ここで $b=(K+cn)/r$、$f=(K-cn)/r$ を用いた。$z=r^2,\lambda=w/r$ と
第二表示の交換にも同じ先頭項の積を使う。定理\ref{thm:rank-fifteen-boundary-disks}で
与えた整数重みによって第 $p$ 位置の第 $j$ 矢を $u^{w_{pj}/2}$ 倍する。
重みの表では $r$ が大きい場合に \texttt{pole-}、小さい場合に \texttt{pole+} を使い、
$w$ が大きい場合に符号 $+1$、小さい場合に $-1$ を使う。
各関係を行ごとに最低次数で割って特殊化する。

\path{computations/verify_rank_fifteen_pencil_infinity.py} は関係の各係数を既約分数にし、
分子・分母それぞれの最低次数の全単項式を $\mathbb Q(v)$ で加える。
その和の非消滅に加え、特殊化に使う全ての先頭係数が剰余でも単元であることを確認する。
$101$ で $v=2$ を使うと二表示と $r,w$ の四選択すべてで、隣接階数6、経路長4,5の
次元35,56、中心文字次元 $(3,2,0,0)$ および全四始点の長さ13までの標準商を得る。
これらの条件は $v$ の非空開集合に延びる。$v\mapsto c^2$ の一般ファイバーの
全てを含むように、有限個の例外値の像を除くことができる。

種数3の二水準では $c=j^2$、$j^4-j^2+1=0$ と取る。
109でこの式は四つの単根に分かれ、$v=j^2+j$ が必要な $L=j$ の両符号を含む。
全四根、二表示、四つの $r,w$ の選択で同じ剰余検算を行った。
一般窓の道次元の下界は重み付き特殊化でも保たれる。道と中心が剰余で一般次元を持つので
中心を持ち上げ、四中心商から全次数の自由性と正則性を得る。
これは重み付き極限そのものにも、その近くで $u$ の付値が正の通常窓にも適用できる。
\end{proof}
